% !TEX program = platex
% 数学ノート集・共通テンプレート v2（積分論2026年度第1部の版面を参考）
% A5判。章内で定義・定理・例・注意等の番号を共通化。
% コンパイル： platex -interaction=nonstopmode math_note_template_v2.tex (2回)
%              dvipdfmx math_note_template_v2.dvi
\documentclass[a5paper,10pt,onecolumn]{jsarticle}
\usepackage[dvipdfmx,top=22mm,bottom=20mm,inner=19mm,outer=18mm,headheight=13pt,headsep=6mm]{geometry}
\usepackage{amsmath,amssymb,amsthm,mathtools}
\usepackage{enumitem,fancyhdr,titlesec}
\usepackage[dvipdfmx,bookmarks=true,colorlinks=false]{hyperref}
\usepackage{pxjahyper}

% ----- 差し替える文書情報 -----
\newcommand{\NoteTitle}{97.120 約数と群作用}
\newcommand{\NoteSeries}{数学ノート集}

% ----- 文字・組版 -----
% 参考PDFの質素な明朝＋Computer Modernの印象を引き継ぎながら，
% A5判で窮屈にならないよう本文の行間を少し確保する．
\linespread{1.13}
\setlength{\parindent}{1zw}
\setlength{\parskip}{0pt}
\setlength{\abovedisplayskip}{6pt plus 2pt minus 2pt}
\setlength{\belowdisplayskip}{6pt plus 2pt minus 2pt}
\setlength{\abovedisplayshortskip}{3pt plus 2pt minus 1pt}
\setlength{\belowdisplayshortskip}{4pt plus 2pt minus 1pt}
\setlist[enumerate]{topsep=3pt,itemsep=1pt,parsep=0pt,leftmargin=2.4zw}
\titleformat{\section}{\normalfont\normalsize\bfseries}{\thesection}{0.8zw}{}
\titlespacing*{\section}{0pt}{2.0ex plus .4ex}{.9ex}
\titleformat{\subsection}{\normalfont\normalsize\bfseries}{\thesubsection}{0.7zw}{}
\titlespacing*{\subsection}{0pt}{1.5ex plus .3ex}{.6ex}
\raggedbottom

% ----- 柱とページ番号：ページ上部、罫線なし -----
\fancypagestyle{mathnote}{%
  \fancyhf{}%
  \fancyhead[L]{\footnotesize \NoteSeries\ （\NoteTitle）}%
  \fancyhead[R]{\footnotesize \thepage}%
  \renewcommand{\headrulewidth}{0pt}%
  \renewcommand{\footrulewidth}{0pt}%
}
\pagestyle{mathnote}

% ----- 節内通し番号の定義・定理群 -----
\newtheoremstyle{mathnotetheorem}
 {4.5pt}{4.5pt}{\normalfont}{}%
 {\bfseries}{}{0.65zw}%
 {\thmname{#1}\thmnumber{ #2}\thmnote{（#3）}}
\theoremstyle{mathnotetheorem}
\newtheorem{theorem}{定理}[section]
\newtheorem{definition}[theorem]{定義}
\newtheorem{lemma}[theorem]{補題}
\newtheorem{proposition}[theorem]{命題}
\newtheorem{corollary}[theorem]{系}
\newtheorem{example}[theorem]{例}
\newtheorem{remark}[theorem]{注意}
\renewcommand{\proofname}{証明}
\renewcommand{\qedsymbol}{$\square$}
\makeatletter
\renewenvironment{proof}[1][\proofname]{%
  \par\pushQED{\qed}\normalfont\topsep4\p@\@plus2\p@\relax
  \trivlist\item[\hskip\labelsep\bfseries #1]\ignorespaces
}{\popQED\endtrivlist\@endpefalse}
\makeatother

% ----- 使うノートで適宜拡張する記号 -----
\newcommand{\R}{\mathbb{R}}
\newcommand{\norm}[1]{\lVert #1\rVert}
\newcommand{\abs}[1]{\lvert #1\rvert}

\begin{document}
\section{問題の設定と主結果}
本ノートでは，約数の対合と有限群の軌道計数を用いて，平方数の条件を満たす整数を数える．
対象読者は初等整数論と群作用を既習とする．ただし必要な軌道計数公式は証明する．

\begin{example}[東京大学2025年度前期・理科数学第4問（問題の要約）]\label{ex:todai}
正整数 $a$ に対し $f_a(n)=n^2+n-a$ とする．
$f_a(n)$ が非負整数の平方となる正整数 $n$ の個数 $N_a$ を考える．
出題では $n\le a$ を示し，$N_a=1$ と $4a+1$ が素数であることの同値性を証明する．
\end{example}

\begin{theorem}[約数個数による完全な解答]\label{thm:main}
$a\ge1$ を整数とし，$M=4a+1$ とする．
$\tau(M)$ を $M$ の正の約数の個数，$\varepsilon(M)$ を $M$ が完全平方数のとき $1$，そうでなければ $0$ とする．
このとき
\begin{equation}\label{eq:main}
 N_a=\frac{\tau(4a+1)+\varepsilon(4a+1)}{2}.
\end{equation}
特に $N_a=1$ であることと $4a+1$ が素数であることは同値である．
\end{theorem}

\section{平方数条件と約数対}
\begin{lemma}[連続する二平方数による評価]\label{lem:bound}
正整数 $a,n$ に対し $n^2+n-a$ が平方数ならば $n\le a$ である．
\end{lemma}
\begin{proof}
$n^2+n-a=m^2$ とする．ここで $m$ は非負整数である．
もし $n>a$ ならば
\[
 n^2<m^2=n^2+(n-a)<n^2+2n+1=(n+1)^2
\]
となる．二つの連続する平方数の間に平方数は存在しないから矛盾である．
よって $n\le a$．さらに $m^2\le n^2$ なので $0\le m\le n$ である．
\end{proof}

\begin{lemma}[二平方差の約数対応]\label{lem:pairs}
$M>1$ を $M\equiv1\pmod 4$ を満たす整数とする．
$M=4a+1$ と書けば，組
\[
 (n,m)\in\mathbb Z_{>0}\times\mathbb Z_{\ge0},\qquad n^2+n-a=m^2
\]
と，正の約数対 $de=M$，$d\le e$ の間には全単射がある．
対応は
\[
 d=2n+1-2m,\quad e=2n+1+2m
\]
で与えられ，逆対応は
\begin{equation}\label{eq:inverse}
 n=\frac{d+e-2}{4},\qquad m=\frac{e-d}{4}
\end{equation}
である．
\end{lemma}
\begin{proof}
組 $(n,m)$ が与えられたとき，補題~\ref{lem:bound} より $0\le m\le n$ であり，$d\ge1$ である．
平方差の因数分解から
\[
 (2n+1-2m)(2n+1+2m)
 =(2n+1)^2-4m^2=4a+1=M
\]
となる．したがって $d,e$ は求める正の約数対である．

逆に，$de=M$，$1\le d\le e$ とする．$M$ が奇数なので $d,e$ は奇数である．
また $de\equiv1\pmod4$ より $d\equiv e\pmod4$ である．
従って式~\eqref{eq:inverse} の右辺は整数であり，$m\ge0$ である．
$d+e\ge 1+M\ge6$ なので $n\ge1$ となる．
さらに
\[
 (2n+1)^2-(2m)^2=de=M=4a+1
\]
から $n^2+n-a=m^2$ を得る．
両操作が互いに逆であることは定義式から直ちに確かめられる．
\end{proof}

\begin{remark}
$M>1$ の仮定は $a\ge1$ に対応する．$M=1$ の場合には $d=e=1$ から $n=0$ となるため，正整数 $n$ の計数では除外する必要がある．
\end{remark}

\section{軌道計数とBurnsideの補題}
\begin{definition}[群作用と軌道]
有限群 $G$ の有限集合 $X$ への作用とは，写像 $G\times X\to X$ で，$1x=x$ および $(gh)x=g(hx)$ を満たすものをいう．
$x\in X$ の軌道を $Gx=\{gx:g\in G\}$ と書き，$g\in G$ の固定点集合を $X^g=\{x\in X:gx=x\}$ とする．
\end{definition}
\begin{lemma}[Burnsideの軌道計数公式]\label{lem:burnside}
有限群 $G$ が有限集合 $X$ に作用するとき，軌道の個数は
\[
 \#(X/G)=\frac{1}{|G|}\sum_{g\in G}|X^g|
\]
である．
\end{lemma}
\begin{proof}
集合 $\Omega=\{(g,x)\in G\times X:gx=x\}$ を二通りに数える．
第一に $|\Omega|=\sum_{g\in G}|X^g|$ である．
第二に各 $x$ の安定化部分群 $G_x$ を用いれば
$|\Omega|=\sum_{x\in X}|G_x|$ である．
軌道 $O=Gx$ 上で $|G_y|=|G|/|O|$ が成立する．
従って一つの軌道からの寄与は $\sum_{y\in O}|G_y|=|G|$ であり，$|\Omega|=|G|\#(X/G)$ を得る．
\end{proof}

\begin{proposition}[約数対は二元群の軌道]\label{prop:orbit}
正整数 $M$ の約数集合 $D(M)=\{d\in\mathbb Z_{>0}:d\mid M\}$ 上で
\[
 \iota(d)=M/d
\]
と定める．このとき $\iota^2=\mathrm{id}$ であるから，$C_2$ が $D(M)$ に作用し，軌道は順序を無視した約数対 $\{d,M/d\}$ と一致する．
軌道の個数は
\[
 \frac{\tau(M)+\varepsilon(M)}2
\]
である．
\end{proposition}
\begin{proof}
$\iota(\iota(d))=M/(M/d)=d$ より対合である．
固定点条件は $d=M/d$，すなわち $d^2=M$ である．
固定点は $M$ が平方数のときただ一つ，そうでなければ存在しない．
補題~\ref{lem:burnside} を $G=C_2$ に適用すれば
\[
 \#(D(M)/C_2)
 =\tfrac12\bigl(|D(M)|+|D(M)^\iota|\bigr)
 =\tfrac12(\tau(M)+\varepsilon(M)).
\]
\end{proof}

\begin{proof}[定理~\ref{thm:main} の証明]
補題~\ref{lem:pairs} により $N_a$ は $4a+1$ の順序を無視した正の約数対の数であり，命題~\ref{prop:orbit} が式~\eqref{eq:main} を与える．
$M=4a+1\ge5$ である．もし $M$ が素数ならば $\tau(M)=2,\varepsilon(M)=0$ なので $N_a=1$．
逆に $M$ が合成数で平方数ならば $\tau(M)\ge3$，平方数でなければ $\tau(M)\ge4$ であるから，いずれも $N_a\ge2$．
従って $N_a=1$ は $M$ の素数性と同値である．
\end{proof}

\section{約数関数による分類}
\begin{proposition}[素因数分解による解の個数]\label{prop:factor}
$4a+1=\prod_{j=1}^r p_j^{e_j}$ を相異なる素数の積に分解する．このとき
\[
 N_a=\frac12\left\{\prod_{j=1}^r(e_j+1)
  +\prod_{j=1}^r\boldsymbol1_{\{e_j\equiv0\pmod2\}}\right\}.
\]
\end{proposition}
\begin{proof}
正の約数は指数 $0\le f_j\le e_j$ の組と一対一に対応するので $\tau(M)=\prod(e_j+1)$ である．
$M$ が平方数であることと，すべての $e_j$ が偶数であることは同値である．
これを定理~\ref{thm:main} に代入する．
\end{proof}

\begin{example}[平方数の特別な扱い]
$a=6$ ならば $M=25$ で，約数対は $(1,25),(5,5)$ の二つである．
式~\eqref{eq:inverse} により対応する組は $(n,m)=(6,6),(2,0)$ であり，$N_6=2$．
平方根 $5$ の固定点が一つ現れることが $\varepsilon(25)=1$ に対応する．
\end{example}

\begin{example}[異なる素数指数型]
$a=11$ ならば $M=45=3^2\cdot5$ である．
$\tau(45)=3\cdot2=6$，平方数でないから $N_{11}=3$．
約数対 $(1,45),(3,15),(5,9)$ に対応する $n$ はそれぞれ $11,4,3$ である．
\end{example}

\begin{corollary}[各個数の出現]
任意の正整数 $k$ に対して，$N_a=k$ を満たす正整数 $a$ が存在する．
\end{corollary}
\begin{proof}
奇素数 $p\equiv1\pmod4$，例えば $p=5$ をとる．
$M=p^{2k-1}$ とすれば $M\equiv1\pmod4$ であり，$a=(M-1)/4$ は正整数である．
$\tau(M)=2k$，$M$ は平方数でないから，定理~\ref{thm:main} により $N_a=k$．
\end{proof}

\section{解の明示と一般化}
\begin{theorem}[解集合の明示]
$a\ge1$，$M=4a+1$ とする．求める正整数 $n$ の集合は
\[
 \left\{\frac{d+M/d-2}{4}:\ d\mid M,\quad 1\le d\le\sqrt M\right\}
\]
である．この集合の各項は異なる．
\end{theorem}
\begin{proof}
集合の記述は補題~\ref{lem:pairs} に従う．
異なる $d_1<d_2\le\sqrt M$ に対して
\[
 d_1+M/d_1-(d_2+M/d_2)
 =(d_2-d_1)\bigl(M/(d_1d_2)-1\bigr)>0
\]
が成立する．最後の不等号は $d_1d_2<M$ による．
よって対応する $n$ は互いに異なる．
\end{proof}

\begin{remark}[一般の差の平方数]
正整数 $M$ の整数による表示 $X^2-Y^2=M$ は，$(X-Y)(X+Y)=M$ への変換で分類できる．
ただし一般には因子が同じ偶奇をもつ条件が必要である．
本ノートでは $X=2n+1$，$Y=2m$ に制限されるため，因子に追加の合同条件が課される．
$M\equiv1\pmod4$ という仮定がこの合同条件を自動的に満たす点が重要である．
\end{remark}

\begin{remark}[証明の依存関係]
本題の独自の論証は，平方数の間隔による評価，因数分解による全単射，有限群の軌道計数の三段階で構成した．
平方数の場合を見落とすと固定点を数えず，式~\eqref{eq:main} の右辺が誤る．
\end{remark}

\section*{参考文献・出典}
\begin{enumerate}[label={[\arabic*]}]
\item 東京大学，2025年度前期日程・理科数学，第4問（問題の出典）．
\item S.~Lang, \textit{Algebra}, revised 3rd ed., Springer, 2002. Group actions and finite groups.
\item H.~Cohen, \textit{A Course in Computational Algebraic Number Theory}, Springer, 1993. Integer factorization and elementary number theory.
\item 「なかけんの数学ノート」，東京大学理系2025年度第4問（問題文の出題条件を照合）．\\
\texttt{https://math.nakaken88.com/problem/tokyo-u-r-2025-4/}
\end{enumerate}

\end{document}